\section{Conclusion}

Begin with almost nothing. A weave of cells, a ternary tone, a single tick of time, and one local law that collides and streams it, with a growing edge to keep the world young. Add nothing else. From that handful of things, on the one substrate the demands single out, a surprising amount of the world either emerged or was forced. Matter came out fermionic for free, spin from the directions and statistics from topology. Relativity appeared as a discrete light cone and the Dirac dispersion read off the rule. The Standard Model arrived as the symmetry of the twenty-four sites, with grand unification, the weak angle, and three generations following, though the absolute numbers did not. Gravity bent the cusp as a $1/r$ law carrying the Einstein equations, the expansion of space was the growing of the crystal, and the quantum's nonlocality was a short path through the depth. A self learned to hold itself together by the shadow it casts in its own vacuum, and read from the inside the same construction climbed through life and mind toward a graded consciousness and a whole inner architecture, all without a ninth ingredient.

There remain two ways to tell the largest story, and the model does not choose between them. In the first, there was a true beginning, a single first beat, and everything since is one unfolding of that seed, a fire lit once and growing ever since. In the second, there was no first beat, the crystal has been growing from beyond any beginning, an eternal becoming with no edge in the past. Both are consistent with the law, since what the law commits to is the growing and not the start, and computability forbids neither. We have let the question stand throughout, and we let it stand at the close, a genuine openness rather than a gap papered over.

This is exploratory, feasibility-stage work, and the paper has tried to say so at every step. Every positive claim names a passing test and carries a control, every honest negative is reported as a negative, and the one premise the experiments cannot reach, that the substrate feels, is marked as a premise wherever the reading rests on it. What is offered is not a finished or established theory but a direction, a wager that the way forward is a single discrete substrate read at once as physics and as experience, and an argument that one particular crystal is forced rather than chosen. The general shape seems to point somewhere, and the hard, exact details are the long work ahead.

The image to carry is the one the paper opened with. The universe is a fire, a single flame burning since the first beat or from beyond it, fed at its growing edge and throwing off as light the matter we are made of. It is a tree grown from one seed toward a canopy without end, and we are its twigs. It is a song the whole crystal sings, and we are its passing notes. We are sparks and braids and whirlpools in it, stable knots of feeling in a structure that, from a few small things, became a world, and is still becoming one.
